\subsection{System overview and pipeline}
\label{ssec:Architecture}
The methodology follows the same order as the implemented workflow. The first
step generates a synthetic, network-aware demand from the SUMO road network
and the Traffic Analysis Zone (TAZ) layer described in
Section~\ref{ssec:stage_1_demand_generation}. This demand is written as
\texttt{trips.xml}: every trip has a departure time, an origin edge and a
destination edge, while the full path is left open. The closure is then
translated into candidate deviation plans. In the Monte Carlo stage, the same
trip demand is routed repeatedly with
\texttt{duarouter},
using randomized edge weights, before each deviation plan is applied and
simulated in SUMO. The last step compares the plans using paired statistics
and a multi-criteria decision rule. The functional organization of this
workflow is summarized in Table~\ref{tab:functional_pipeline}, while
Figure~\ref{fig:pipeline} shows the same pipeline as a data flow from inputs to
final evaluation files.

\subsubsection{Concepts and contribution boundary}

The methods chapter combines standard transport concepts, external SUMO tools
and code developed for this thesis. Table~\ref{tab:method_concepts} defines
the main terms before the individual stages are described.

\begin{longtable}{@{}p{0.22\linewidth}>{\raggedright\arraybackslash}p{0.72\linewidth}@{}}

\toprule
Term & Meaning in this thesis \\
\midrule
\endfirsthead

\toprule
Term & Meaning in this thesis \\
\midrule
\endhead

\endfoot

\bottomrule

\caption{Key concepts used in the methodology and the boundary between reused
methods/tools and thesis-specific implementation choices.}
\label{tab:method_concepts}
\endlastfoot

SUMO & \emph{Simulation of Urban MObility}, the external microscopic traffic simulator used to execute the traffic runs. SUMO is used as a tool, not developed in this thesis. \\

\texttt{duarouter}\footnotemark & SUMO's command-line router. It converts unrouted trips into explicit route files, this thesis uses its randomized edge-weight option to generate route-choice variation. \\

TAZ & \emph{Traffic Analysis Zone}, a geographic zone used to aggregate origins and destinations before individual trips are instantiated on SUMO road edges. \\

OD matrix & \emph{Origin-Destination} matrix. Entry $(i,j)$ gives the number or share of trips from origin zone $i$ to destination zone $j$. \\

Gravity model \cite{gravity_model} & A standard demand-allocation model from transport planning. The thesis contribution is the synthetic implementation: zone-mass proxy, sampled TAZ-to-TAZ impedance and integer hourly OD matrices. \\

Deviation plan & A candidate detour path around the closed edge or edge sequence. The plan objects and route-rewriting policy are implemented in this thesis. \\

Baseline & The no-detour, no-closure condition for the same demand and Monte Carlo seed. Each candidate plan is compared with its paired baseline run. \\

Monte Carlo run & One repeated simulation under a fixed trip set but a different routing seed and randomized edge-cost realization. \\

BFS & \emph{Breadth-first search}, a standard graph search used here to find upstream candidate detour origins by hop depth. \\

MCDA & \emph{Multi-criteria decision analysis}, the final ranking layer that combines time-loss delay, mean travel time and impacted-user delay after feasibility and Pareto screening. \\

\end{longtable}
\footnotetext{\url{https://sumo.dlr.de/docs/CLI.html}}

The core contribution is therefore not a new traffic simulator, a new gravity
model or a new shortest-path algorithm. It is the integration of these
standard components into a reproducible pipeline that generates synthetic
demand, computes candidate detours, applies those detours to route files and
evaluates them across paired stochastic SUMO simulations.

Table~\ref{tab:functional_pipeline} identifies the four implemented modules,
their role in the workflow and the artifacts they write for the following
stage. The table should be read as the functional view of the system: it
explains what each module is responsible for before the following subsections
describe the internal methods in more detail.

\begin{table}[htbp]
\centering
\begin{tabularx}{\linewidth}{@{}l>{\raggedright\arraybackslash}X>{\raggedright\arraybackslash}X@{}}
\toprule
Module & Main role & Main outputs \\
\midrule
\texttt{scenario\_generator} & Generates synthetic OD demand and unrouted trip records & \texttt{trips.xml}, OD matrices \\
\texttt{deviation\_plan\_maker} & Converts the network into a graph and calculates the deviation plan(s) & JSON of candidate deviation plans \\
\texttt{monte\_carlo\_simulation} & Re-routes trip demand and simulates baseline and detour conditions over multiple runs & \texttt{run\_\#} folders, \texttt{summary.json} \\
\texttt{evaluator} & Paired statistical analysis and multi-criteria ranking & \texttt{evaluation.json}, \texttt{evaluation.md} \\
\bottomrule
\end{tabularx}
\caption{Functional organization of the pipeline.}
\label{tab:functional_pipeline}
\end{table}

Figure~\ref{fig:pipeline} complements the table with the execution order. The
dashed box contains the three inputs shared by the workflow, the vertical
arrows show the order in which modules are executed and the side boxes show
the main files produced along the way.

\begin{figure}[H]
\centering
\begin{tikzpicture}[
  module/.style={
    draw, rounded corners=5pt, fill=blue!7, thick,
    minimum width=5.8cm, minimum height=1.0cm,
    align=center, font=\small\ttfamily
  },
  artifact/.style={
    draw, rounded corners=3pt, fill=gray!8,
    minimum width=4.6cm, minimum height=0.75cm,
    align=center, font=\scriptsize\ttfamily
  },
  inputbox/.style={
    draw, dashed, rounded corners=3pt, fill=orange!7,
    minimum width=5.8cm, minimum height=0.85cm,
    align=center, font=\small
  },
  arr/.style={-{Stealth[length=6pt]}, thick},
  side/.style={-{Stealth[length=5pt]}, thick, gray!70},
  node distance=0.75cm
]
  \node[inputbox] (inp)
    {\texttt{.net.xml} \enspace \texttt{taz.xml} \enspace \texttt{config.toml}};

  \node[module, below=0.85cm of inp]  (sg)  {scenario\_generator};
  \node[module, below=0.75cm of sg]   (dpm) {deviation\_plan\_maker};
  \node[module, below=0.75cm of dpm]  (mcs) {monte\_carlo\_simulation};
  \node[module, below=0.75cm of mcs]  (ev)  {evaluator};

  \node[artifact, right=1.6cm of sg]  (sg_out)  {\texttt{trips.xml}, OD matrices};
  \node[artifact, right=1.6cm of dpm] (dpm_out) {\texttt{plans.json}};
  \node[artifact, right=1.6cm of mcs] (mcs_out) {\texttt{run\_\#/}, \texttt{summary.json}};
  \node[artifact, right=1.6cm of ev]  (ev_out)  {\texttt{evaluation.json}};

  \draw[arr] (inp) -- (sg);
  \draw[arr] (sg)  -- (dpm);
  \draw[arr] (dpm) -- (mcs);
  \draw[arr] (mcs) -- (ev);

  \draw[side] (sg)  -- (sg_out);
  \draw[side] (dpm) -- (dpm_out);
  \draw[side] (mcs) -- (mcs_out);
  \draw[side] (ev)  -- (ev_out);
\end{tikzpicture}
\caption{Four-module pipeline. Inputs (dashed box) flow top-to-bottom through the four
modules. The scenario generator writes unrouted trips, explicit route files are produced
later inside each Monte Carlo iteration by duarouter.}
\label{fig:pipeline}
\end{figure}
